%\documentclass[tikz,border=5pt]{standalone}
\documentclass[12pt,a4paper]{article}

% Packages
\usepackage[margin=2.5cm]{geometry}
\usepackage{amsmath, amssymb}
\usepackage{graphicx}
\usepackage{xcolor}
\usepackage{listings}
\usepackage{booktabs}
\usepackage{fancyhdr}
\usepackage{titlesec}
\usepackage{hyperref}
\usepackage{enumitem}
\usepackage{tcolorbox}
\usepackage{array}
\usepackage{tabularx}
\usepackage{caption}
\usepackage{subcaption}
\usepackage{float}
\usepackage{mdframed}
\usepackage{tikz}
\usepackage{bm}
\usepackage{varwidth}

% Figure helper: include the figure when available; otherwise show a clear placeholder.
\newcommand{\figureplaceholder}[2]{%
    \begin{figure}[H]
    \centering
    \IfFileExists{#2}{%
        \includegraphics[width=0.88\textwidth,height=0.42\textheight,keepaspectratio]{#2}%
    }{%
        \fbox{\parbox[c][3cm][c]{0.82\textwidth}{\centering\small Figure not found: \texttt{\detokenize{#2}}}}%
    }%
    \caption{#1}
    \end{figure}
}

% Equation box style
\newsavebox{\equationboxsave}
\newenvironment{equationbox}{
    \begin{center}
    \begin{lrbox}{\equationboxsave}
    \begin{varwidth}{\textwidth}
}{
    \end{varwidth}
    \end{lrbox}
    \begin{tcolorbox}[
        colback=blue!5,
        colframe=blue!45!black,
        boxrule=1pt,
        arc=8pt,
        left=10pt,
        right=10pt,
        top=4pt,
        bottom=4pt,
        width=\dimexpr\wd\equationboxsave+22pt\relax,
        before skip=8pt,
        after skip=8pt
    ]
    \usebox{\equationboxsave}
    \end{tcolorbox}
    \end{center}
}

% ============================================================
% Python Code Formatting
% ============================================================

% Define Python code style
\definecolor{codegreen}{rgb}{0.0,0.45,0.25}
\definecolor{codeblue}{rgb}{0.0,0.25,0.65}
\definecolor{codepurple}{rgb}{0.50,0.0,0.50}
\definecolor{codered}{rgb}{0.65,0.10,0.10}
\definecolor{codegray}{rgb}{0.45,0.45,0.45}
\definecolor{codebackground}{rgb}{0.97,0.97,0.97}
\definecolor{darkblue}{RGB}{0,51,102}
\definecolor{navyblue}{RGB}{0,0,128}

\lstdefinestyle{pythonstyle}{
    language=Python,
    backgroundcolor=\color{codebackground},
    basicstyle=\ttfamily\small,
    keywordstyle=\color{codeblue}\bfseries,
    commentstyle=\color{codegreen},
    stringstyle=\color{codered},
    identifierstyle=\color{black},
    numberstyle=\tiny\color{codegray},
    numbers=left,
    numbersep=8pt,
    frame=single,
    rulecolor=\color{gray!40},
    breaklines=true,
    breakatwhitespace=false,
    showstringspaces=false,
    tabsize=4,
    keepspaces=true,
    columns=fullflexible
}

\lstset{
    language=Verilog,
    basicstyle=\ttfamily\small,
    keywordstyle=\color{blue},
    commentstyle=\color{green!60!black},
    stringstyle=\color{red},
    breaklines=true,
    frame=single,
    numbers=left,
    numberstyle=\tiny,
    style=pythonstyle,
}

\tikzstyle{block} = [rectangle, draw, minimum width=2.5cm, minimum height=0.9cm, text centered, font=\small]
\tikzstyle{arrow} = [->, thick]

% =========================
% Page Setup
% =========================
\geometry{
    top=1in,
    bottom=1in,
    left=1in,
    right=1in
}


\hypersetup{
    colorlinks=true,
    linkcolor=black,
    urlcolor=blue,
    citecolor=black
}


% Header/Footer
\setlength{\headheight}{25pt}
\pagestyle{fancy}
\fancyhf{}
\fancyhead[L]{\small\textcolor{black}{\textbf{EN3150 }}}
\fancyhead[R]{\small\textcolor{black}{Pattern Recognition}}
\fancyfoot[C]{\small\textcolor{black}{\thepage}}
\renewcommand{\headrulewidth}{1.5pt}
\renewcommand{\headrule}{\hbox to\headwidth{\color{black}\leaders\hrule height \headrulewidth\hfill}}

% Section formatting
\titleformat{\section}
  {\large\bfseries\color{darkblue}}
  {\thesection}{1em}{}
  [\vspace{-4pt}\textcolor{darkblue}{\hrule height 0.8pt}\vspace{2pt}]
\titleformat{\subsection}
  {\normalsize\bfseries\color{navyblue}}
  {\thesubsection}{1em}{}
\titlespacing{\section}{0pt}{16pt}{8pt}
\titlespacing{\subsection}{0pt}{10pt}{4pt}
\usetikzlibrary{shapes.geometric, arrows, positioning, fit}
\setcounter{tocdepth}{3}

% =========================
% Document
% =========================
\begin{document}

% =========================
% Title Page
% =========================
\begin{titlepage}
    \centering
    \vspace*{1cm}

    % Institution
    {\LARGE\textbf{\textcolor{codered}{University of Moratuwa}}\par}
    {\normalsize{\textcolor{black}{Department of Electronic and Telecommunication Engineering}}\par}
    \vspace{0.3cm}
    \textcolor{black}{\rule{0.6\textwidth}{1.5pt}}
    \vspace{1.2cm}

    % Course badge
    \begin{figure}[H]
    \centering
    \IfFileExists{uom_logo.png}{\includegraphics[width=7cm,height=7cm]{uom_logo.png}}{\fbox{\parbox[c][5cm][c]{7cm}{\centering University of Moratuwa}}}
    \end{figure}

    \vspace{1.2cm}

    {\LARGE\bfseries\textcolor{darkblue}{Resource-Constrained CNN for Edge Image Classification}\par}
    \vspace{0.4cm}
    {\Large {EN3150 Pattern Recognition}\par}

    \vspace{1.5cm}

    % Details table
    \begin{table}[h]
    \centering
    \begin{tabular}{ll}
    \textbf{Group Name: } & \textbf{Rubix} \\
    \\
    \textbf{230264R} & Janurajah J. \\
    \textbf{230585C} & Saruka U.\\
    \textbf{230638T} & Thayarathan V.\\
    \\
    \textbf{Submitted on} & October 19, 2026 \\
    \end{tabular}
\end{table}
  
\end{titlepage}

% =========================
% Table of Contents
% =========================
\tableofcontents

\newpage
\pagenumbering{arabic}

% -------------------- 1. INTRODUCTION --------------------
\section{Introduction}
Convolutional neural networks (CNNs) are widely used for image classification, but the memory and computational requirements of large models can make deployment difficult on embedded devices and edge nodes. This assignment investigates the trade-off between classification performance and resource efficiency by developing two custom CNNs and comparing them with two pretrained lightweight architectures.

Fashion-MNIST is used as the image-classification dataset. Model A uses conventional two-dimensional convolutions, while Model B uses depthwise separable convolutions and is constrained to no more than 100,000 trainable parameters. MobileNetV2 and EfficientNetB0 are fine-tuned as pretrained reference models.

The objectives are to prepare low-resolution data with a 70/15/15 split; design and document two custom CNN architectures; compare Adam, SGD and SGD with Momentum for Model B; evaluate all models using classification metrics and resource measures; and discuss accuracy, memory footprint and computational cost.

% -------------------- 2. DATA --------------------
\section{Dataset and Preprocessing}
\subsection{Dataset description}
Fashion-MNIST contains 70,000 grayscale images of clothing items across ten classes. Each original image is 28 $\times$ 28 pixels. Table~\ref{tab:classes} lists the labels.
\begin{table}[H]\centering
\caption{Fashion-MNIST class labels}\label{tab:classes}
\begin{tabular}{cl}\toprule Label & Class name\\\midrule
0 & T-shirt/top\\1 & Trouser\\2 & Pullover\\3 & Dress\\4 & Coat\\5 & Sandal\\6 & Shirt\\7 & Sneaker\\8 & Bag\\9 & Ankle boot\\\bottomrule
\end{tabular}\end{table}

\subsection{Data split}
The data were split into training, validation and test subsets using a 70/15/15 ratio. The processed dataset used in all five notebooks contains 70,000 examples split into 49,000 training images, 10,500 validation images and 10,500 test images. The proportions are therefore 70\%, 15\% and 15\%.
\begin{table}[H]\centering
\caption{Dataset split used in the experiments}\label{tab:split}
\begin{tabular}{lrr}\toprule Subset & Percentage & Number of images\\\midrule
Training & 70\% & 49,000\\
Validation & 15\% & 10,500\\
Testing & 15\% & 10,500\\\midrule
Total & 100\% & 70,000\\\bottomrule
\end{tabular}\end{table}

\subsection{Preprocessing pipeline}
For the custom CNNs, images are represented as single-channel grayscale inputs at 28 $\times$ 28 pixels, normalized to the range $[0,1]$. MobileNetV2 and EfficientNetB0 use images resized to 64 $\times$ 64 pixels and converted to three channels. The same underlying training, validation and test examples are used for all models.

For MobileNetV2, the notebook applies the Keras MobileNetV2 preprocessing function. The EfficientNetB0 Keras application includes its own rescaling layer, so the notebook passes the resized float32 image values without applying the MobileNetV2 transformation. The pipelines are therefore model-specific.

\figureplaceholder{Representative Fashion-MNIST images with class labels.}{fig/samples.png}

% -------------------- 3. ARCHITECTURES --------------------
\newpage
\section{Custom CNN Architecture Design}
\subsection{Model A: Standard CNN}
Model A is the baseline CNN. Conventional convolutional layers extract spatial features, and max-pooling layers reduce feature-map dimensions. The resulting representation is flattened and passed to a dense hidden layer and a ten-class softmax output.

\begin{table}[H]\centering
\caption{Model A architecture and parameter counts from the executed model summary.}\label{tab:modela}
\begin{tabularx}{\textwidth}{l X r}\toprule
Layer & Configuration & Parameters\\\midrule
Input & $28\times28\times1$ & 0\\
Conv2D & 16 filters, $3\times3$, ReLU & 160\\
MaxPooling2D & $2\times2$ & 0\\
Conv2D & 32 filters, $3\times3$, ReLU & 4,640\\
MaxPooling2D & $2\times2$ & 0\\
Conv2D & 64 filters, $3\times3$, ReLU & 18,496\\
MaxPooling2D & $2\times2$ & 0\\
Flatten & Flatten feature maps & 0\\
Dense & 64 units, ReLU & 36,928\\
Dense & 10 units, softmax &  650 \\\midrule
Total & & 60,874 \\\bottomrule
\end{tabularx}\end{table}

For a conventional convolution with biases, the parameter count is
\begin{equation}\label{eq:convparams}
P_{\mathrm{conv}}=(K_hK_wC_{\mathrm{in}}+1)C_{\mathrm{out}},
\end{equation}
where $K_h$ and $K_w$ are kernel dimensions, $C_{\mathrm{in}}$ is the number of input channels and $C_{\mathrm{out}}$ is the number of filters. For a dense layer,
\begin{equation}\label{eq:denseparams}
P_{\mathrm{dense}}=(n_{\mathrm{in}}+1)n_{\mathrm{out}}.
\end{equation}

%\figureplaceholder{Model A architecture diagram or screenshot of its model summary. Ensure it matches the executed code.}{02_Model_A_Standard_CNN_cell10_fig1.png}

\subsection{Model B: Lightweight CNN}
Model B replaces conventional convolutions with depthwise separable convolutions. A depthwise convolution applies a spatial filter independently to each input channel; a pointwise $1\times1$ convolution combines the channel-wise features. This factorization can reduce parameter count and multiply-accumulate operations relative to a conventional convolution with comparable channel dimensions.

For a depthwise separable convolution with a $K_h\times K_w$ kernel, $C_{\mathrm{in}}$ input channels and $C_{\mathrm{out}}$ output channels, with one bias per output channel,
\begin{equation}\label{eq:sepparams}
P_{\mathrm{sep}}=K_hK_wC_{\mathrm{in}}+C_{\mathrm{in}}C_{\mathrm{out}}+C_{\mathrm{out}}.
\end{equation}

\begin{table}[H]
\centering
\caption{Model B lightweight CNN architecture, grouped by functional block}
\label{tab:modelb}
\small
\begin{tabularx}{\textwidth}{l X r}
\toprule
Layer/block & Configuration & Parameters\\
\midrule
Input & $28\times28\times1$ & 0\\
Stem convolution & $3\times3$, 16 channels, no bias; batch normalization and ReLU & 208\\
Max pooling & $2\times2$ & 0\\
Depthwise-separable block 1 & Depthwise $3\times3$ (16 channels), pointwise $1\times1$ to 24 channels; batch normalization and ReLU & 688\\
Max pooling & $2\times2$ & 0\\
Depthwise-separable block 2 & Depthwise $3\times3$ (24 channels), pointwise $1\times1$ to 32 channels; batch normalization and ReLU & 1,208\\
Max pooling & $2\times2$ & 0\\
Depthwise-separable block 3 & Depthwise $3\times3$ (32 channels), pointwise $1\times1$ to 48 channels; batch normalization and ReLU & 2,144\\
Global average pooling and dropout & Dropout rate 0.20 & 0\\
Dense output & 10 units, softmax & 490\\
\midrule
Total & 4,354 trainable and 384 non-trainable parameters & 4,738\\
\bottomrule
\end{tabularx}
\end{table}


\subsection{Activation functions and hardware considerations}
ReLU is used in hidden layers because $f(x)=\max(0,x)$ is computationally simple and avoids the exponential calculations used by sigmoid and tanh. It is widely supported by deep-learning runtimes. The output layer uses softmax to convert class logits into a probability distribution over ten classes. Softmax requires exponentials but is appropriate for normalized multi-class probabilities; its cost should be considered separately from the main convolutional computation.

% -------------------- 4. OPTIMIZERS --------------------
\newpage
\section{Optimizer Selection and Training Methodology}
\subsection{Optimizer comparison}
Model B's optimizer selection compares Adam, standard stochastic gradient descent (SGD), and SGD with Momentum. Candidates should use the same data splits, architecture, batch size and comparable training settings. Select the optimizer using validation metrics, not test-set performance.

\begin{table}[H]\centering
\caption{Optimizer comparison using the recorded validation results.}\label{tab:optimizers}
\small
\setlength{\tabcolsep}{3pt}
\begin{tabularx}{\textwidth}{>{\raggedright\arraybackslash}X
>{\centering\arraybackslash}p{1.8cm}
>{\centering\arraybackslash}p{3.2cm}
>{\centering\arraybackslash}p{2.1cm}
>{\centering\arraybackslash}p{2cm}}\toprule
Optimizer & \shortstack{Learning\\rate} & \shortstack{Final\\validation\\accuracy (\%)} & \shortstack{Final\\validation\\loss} & \shortstack{Time\\(s/epoch)}\\\midrule
Adam & 0.001 & 87.09 & 0.3533 & 6.90\\
SGD & 0.01 & 84.88 & 0.4237 & 6.87\\
SGD with Momentum & 0.01 & 86.50 & 0.3700 & 7.31\\\bottomrule
\end{tabularx}\end{table}

SGD updates parameters using the current gradient. Momentum incorporates a contribution from previous gradients, which can smooth updates and accelerate progress in consistent directions. Adam adapts update scales using estimates of first and second gradient moments. The methods may differ in convergence speed, stability and validation performance.

\textbf{Selected optimizer and learning rate:} Adam, learning rate $10^{-3}$.\\
\textbf{Reason for selection:} Adam achieved the highest final validation accuracy (87.09\%) and the lowest final validation loss (0.3533) in the recorded comparison. SGD with Momentum improved validation accuracy over plain SGD (86.50\% versus 84.88\%), but that run completed only 11 epochs, so its endpoint is not directly comparable with the 20-epoch runs.
\figureplaceholder{Optimizer comparison plot showing validation loss.}{fig/model_b_optimizer.png}

\subsection{Training configuration}
Both custom CNNs completed the required minimum of 20 epochs. The following settings are taken from the executed notebooks.
\begin{table}[H]\centering
\caption{Training configuration}\label{tab:training}
\small
\setlength{\tabcolsep}{4pt}
\begin{tabularx}{\textwidth}{>{\raggedright\arraybackslash}p{3.1cm}
    >{\raggedright\arraybackslash}X
    >{\raggedright\arraybackslash}X}\toprule
Setting & Model A & Model B\\\midrule
Loss function & Sparse categorical cross-entropy & Sparse categorical cross-entropy\\
Optimizer & Adam & Adam\\
Learning rate & 0.001 & 0.001\\
Batch size & 128 & 128\\
Completed epochs & 20 & 25\\
Checkpointing / early stopping & Best validation checkpoint & Best validation checkpoint; LR reduction\\\bottomrule
\end{tabularx}\end{table}

Each pretrained model completed five epochs with the backbone frozen and five fine-tuning epochs. The timing results below are averages over these ten epochs; because notebook execution environments may differ, compare their timings cautiously.

% -------------------- 5. RESULTS --------------------
\newpage
\section{Experimental Results}
\subsection{Model A learning curves}
\figureplaceholder{Model A training and validation loss and accuracy across epochs.}{fig/training_curves_a.png}

\subsection{Model B learning curves}
\figureplaceholder{Model B training and validation loss and accuracy across epochs.}{fig/training_curves_b.png}

\subsection{Learning Curve Analysis}

Model A shows signs of overfitting, as its training loss decreases
while its validation loss increases in later epochs. In contrast,
Model B maintains improving performance with a smaller gap between
training and validation results. EfficientNetB0 also shows steady
improvement, while MobileNetV2's validation performance begins to
level off. Overall, Model B offers a good balance between accuracy
and model complexity.

\subsection{Test-set classification performance}
All four models were evaluated on the same held-out test set of 10,500 images. The metrics below are taken from the saved notebook outputs.
\begin{table}[H]\centering
\caption{Test-set classification metrics}\label{tab:metrics}
\resizebox{\textwidth}{!}{\begin{tabular}{lrrrrr}\toprule
Model & Accuracy (\%) & Weighted precision & Weighted recall & Macro precision & Macro recall\\\midrule
Model A & 89.58 & 0.8963 & 0.8958 & 0.8963 & 0.8958\\
Model B & 89.10 & 0.8902 & 0.8910 & 0.8902 & 0.8910\\
MobileNetV2 & 88.23 & 0.8815 & 0.8823 & 0.8815 & 0.8823\\
EfficientNetB0 & 89.29 & 0.8921 & 0.8929 & 0.8921 & 0.8929\\\bottomrule
\end{tabular}}\end{table}

Precision measures how many predictions assigned to a class are correct; recall measures how many true examples of a class are retrieved. Weighted averages account for class support, while macro averages give each class equal weight. The test split contains exactly 1,050 examples from each class, so weighted and macro precision/recall are equal (up to rounding) for these results.

\subsection{Confusion matrices}
\figureplaceholder{Model A test-set confusion matrix.}{fig/a_matrix.png}
\figureplaceholder{Model B test-set confusion matrix.}{fig/b_matrix.png}
\figureplaceholder{MobileNetV2 test-set confusion matrix.}{fig/mobilenetv2_matrix.png}
\figureplaceholder{EfficientNetB0 test-set confusion matrix.}{fig/efficientnetb0_matrix.png}

\subsection{Pretrained model learning curves}
\figureplaceholder{MobileNetV2 training and validation loss curves.}{fig/training_curves_mobilenet.png}
\figureplaceholder{MobileNetV2 training and validation accuracy curves.}{fig/training_curves_mobilenetv2.png}
\figureplaceholder{EfficientNetB0 training and validation loss curves.}{fig/training_curves_efficientnet.png}
\figureplaceholder{EfficientNetB0 training and validation accuracy curves.}{fig/training_curves_efficientnetb0.png}

% -------------------- 6. COMPARISON --------------------
\newpage
\section{Resource Efficiency and Model Comparison}
\subsection{Custom model comparison}
\begin{table}[H]\centering
\caption{Required comparison of the two custom CNNs}\label{tab:customcomparison}
\begin{tabular}{lrr}\toprule Metric & Model A & Model B\\\midrule
Total parameters & 60,874 & 4,738\\
Estimated/saved model size (KB) & 757.90 & 18.51\\
Average training time per epoch (s) & 17.17 & 7.52\\
Test accuracy (\%) & 89.58 & 89.10\\\bottomrule
\end{tabular}\end{table}

\noindent\textbf{Model size comparison:}
Model A's size (757.90 KB) is measured from its saved
\texttt{.keras} file, whereas Model B's size (18.51 KB)
is estimated from its float32 parameter footprint.
Therefore, these values are not directly comparable.
For a fair comparison, both models should be saved and
measured using the same method.

Model B achieved 89.10\% test accuracy, close to Model A's
89.58\%. However, Model B uses only 4,738 parameters,
compared with 60,874 for Model A, reducing the parameter
count by approximately 92.22\%.

Model B also had a shorter measured training time per epoch
(7.52 seconds compared with 17.17 seconds for Model A).
These results show that a lightweight CNN can achieve
competitive accuracy with far fewer parameters, making it
a promising choice for resource-constrained devices.

However, the model-size figures are not directly comparable:
Model A's size is measured from its saved file, while Model B's
size is an estimated parameter footprint. Therefore, further
testing using the same measurement method is needed for a
fair comparison.

\subsection{Comparison with lightweight pretrained networks}
\begin{table}[H]\centering
\caption{Comparison of Model B with pretrained lightweight models}\label{tab:sotacomparison}
\resizebox{\textwidth}{!}{\begin{tabular}{lrrr}\toprule
Metric & Model B & MobileNetV2 & EfficientNetB0\\\midrule
Total parameters & 4,738 & 2,270,794 & 4,062,381\\
Trainable parameters after fine-tuning & 4,354 & 1,207,690 & 1,355,578\\
Weight size / saved file (MB) & 0.018$^*$ & 18.447 & 26.658\\
Average training time per epoch (s) & 7.52 & 40.39 & 61.17\\
Test accuracy (\%) & 89.10 & 88.23 & 89.29\\
Weighted precision & 0.8902 & 0.8815 & 0.8921\\
Weighted recall & 0.8910 & 0.8823 & 0.8929\\
Input resolution & $28\times28\times1$ & $64\times64\times3$ & $64\times64\times3$\\\bottomrule
\end{tabular}}\end{table}


\subsection{Accuracy, memory and computation trade-offs}

\begin{enumerate}[leftmargin=1.5em]

    \item \textbf{Highest test accuracy:}
    Model A achieved the highest test accuracy of $89.58\%$,
    outperforming Model B ($89.10\%$) by $0.48$ percentage points.

    \item \textbf{Model size:}
    Model B has only $4{,}738$ parameters, compared with
    $2{,}270{,}794$ for MobileNetV2 and $4{,}062{,}381$ for
    EfficientNetB0. Its estimated float32 weight footprint is only
    18.51 KB. The notebook-reported saved sizes for MobileNetV2 and
    EfficientNetB0 are 18.447 MB and 26.658 MB, respectively.

    \item \textbf{Effect of pretraining:}
    ImageNet pretraining did not improve accuracy over Model A.
    MobileNetV2 achieved $88.23\%$, while EfficientNetB0 achieved
    $89.29\%$. However, EfficientNetB0 performed slightly better
    than Model B.

    \item \textbf{Training time:}
    Model B had the shortest measured custom-model epoch time
    (7.52 s). MobileNetV2 and EfficientNetB0 averaged 40.39 s and
    61.17 s per epoch in their recorded runs. The pretrained-model
    notebooks ran in a different environment, so these timings are
    indicative rather than a controlled hardware comparison.
    Training time alone does not determine inference speed or energy
    use, which depend on the hardware and implementation.

    \item \textbf{Model A versus Model B:}
    Model B achieved nearly the same accuracy as Model A while
    reducing the parameter count by approximately $92.22\%$.
    This makes it a useful option when memory is limited, with
    only a small reduction in test accuracy.

    \item \textbf{Overall comparison:}
    Model B offers a compact and efficient architecture, while
    Model A provides the highest test accuracy. Pretrained models
    have substantially more parameters but can benefit from
    previously learned features. The best choice depends on the
    application's accuracy and resource requirements.

\end{enumerate}

The experiments show that a smaller model can achieve competitive accuracy without requiring a large pretrained architecture. The custom-model timings were recorded by the notebooks using elapsed wall-clock time around model fitting. The CPU model and RAM were not recorded, and the pretrained notebooks ran in a separate environment; therefore, absolute timing comparisons should be treated cautiously.

% -------------------- 7. LIMITATIONS --------------------
\newpage
\section{Limitations}
Fashion-MNIST is a low-resolution dataset with ten classes, so results may not generalize to more complex real-world imagery. The pretrained networks receive resized three-channel inputs, whereas custom CNNs can operate directly on grayscale images; this difference should be considered when interpreting computation and input-processing overhead.

Training time depends on CPU/GPU hardware, runtime configuration, batch size, software versions and whether data loading is included. Saved model size depends on serialization format. Parameter count, file size and training time do not alone measure on-device inference latency, peak RAM use, energy consumption or multiply-accumulate operations.

% -------------------- 8. CONCLUSION --------------------
\section{Conclusion}

This assignment compared two custom CNN architectures with two pretrained image-classification networks on Fashion-MNIST. Model A provided a conventional CNN baseline, while Model B used depthwise separable convolutions to reduce the parameter count.

Model A achieved the highest test accuracy of 89.58\%, followed by EfficientNetB0 at 89.29\%, Model B at 89.10\%, and MobileNetV2 at 88.23\%. With only 4,738 parameters, Model B achieved competitive accuracy while using far fewer parameters than the pretrained models. Overall, the results show that a carefully designed lightweight CNN can achieve comparable performance with much lower model complexity, making it a practical choice when memory and computational resources are limited.

% -------------------- 9. REPRODUCIBILITY --------------------
\section{Reproducibility and GitHub Repository}
\textbf{GitHub repository:} \url{https://github.com/Saruka03/EN3150-CNN-Classification}\\
The repository is available at the URL above. For reproducibility, it should include the five Jupyter notebooks, the processed dataset or instructions to recreate it, model-output assets, and a README describing dependencies and execution order.

\textbf{Hardware used:} Notebook logs report no available GPU for the Model B run; CPU model and RAM were not recorded in the exported outputs.\\
\textbf{Software versions:} TensorFlow 2.21.0 is recorded in the Model B notebook; other package versions were not consistently recorded.

% -------------------- REFERENCES --------------------
\newpage
\section{References}
\begin{enumerate}[leftmargin=1.5em]
\item H. Xiao, K. Rasul and R. Vollgraf, ``Fashion-MNIST: a Novel Image Dataset for Benchmarking Machine Learning Algorithms,'' 2017. Available: \url{https://github.com/zalandoresearch/fashion-mnist}.
\item A. G. Howard et al., ``MobileNets: Efficient Convolutional Neural Networks for Mobile Vision Applications,'' arXiv:1704.04861, 2017. Available: \url{https://arxiv.org/abs/1704.04861}.
\item M. Sandler et al., ``MobileNetV2: Inverted Residuals and Linear Bottlenecks,'' CVPR, 2018. Available: \url{https://arxiv.org/abs/1801.04381}.
\item M. Tan and Q. V. Le, ``EfficientNet: Rethinking Model Scaling for Convolutional Neural Networks,'' ICML, 2019. Available: \url{https://arxiv.org/abs/1905.11946}.
\item TensorFlow, ``Transfer learning and fine-tuning.'' Available: \url{https://www.tensorflow.org/guide/keras/transfer_learning}.
\item TensorFlow, ``Keras Applications.'' Available: \url{https://www.tensorflow.org/api_docs/python/tf/keras/applications}.
\item EN3150 Pattern Recognition, ``Assignment 03: Resource-Constrained CNN for Edge Image Classification,'' course assignment brief, University of Moratuwa, 2026.
\end{enumerate}


\end{document}